\chapter{Аудит trace-селекторов superconnection curvature}
\label{ch:v10-h15-r2-superconnection-trace-selector-audit}

Градуированная поляризация суперкривизны дала точный оператор $Q$, но её
trace metric неопределённа. Проверим, может ли другой канонический селектор
сохранить endpoint-разность внутри положительного действия.

Два независимых endpoint-канала образуют матрицу
\begin{equation}
 C=(T,-B),\qquad C^TC=\operatorname{diag}(72,64).
\end{equation}
Уравнение $Cw=Q=T+B$ имеет единственное решение
\begin{equation}
 \boxed{w=(1,-1)}.
\end{equation}
Следовательно, никакой положительный диагональный endpoint-trace с
неотрицательными весами не может дать $Q$. Ordinary trace имеет веса
$(1,1)$ и возвращает $T-B$; raw supertrace возвращает $Q$, но его weight
metric имеет inertia $(1_-,0_0,1_+)$.

Остаётся структурно иной обход: сначала выделить relative curvature block,
а затем взять его положительную норму. Канонический проектор на разностную
линию равен
\begin{equation}
 P_{\rm rel}=\frac12
 \begin{pmatrix}1&-1\\-1&1\end{pmatrix},
 \qquad P_{\rm rel}^2=P_{\rm rel},
 \qquad \rank P_{\rm rel}=1.
\end{equation}
Его Gram form положительно полуопределена с inertia $(0_-,1_0,1_+)$. Это
показывает, что positivity и exact endpoint difference совместимы после
проекции на отдельную curvature-компоненту. Но унаследованный relative
selector текущего представленного calculus имеет rank $0$.

Двенадцать trace-механизмов проверены по шести критериям: exact $Q$,
positivity полного действия, gauge/Real invariance, каноническая
нормировка, совместимость с локальной степенью и inherited selector.
Матрица имеет rank $6$, проходов нет. Ordinary trace получает $5/6$ и
проваливает exact $Q$; raw supertrace получает $5/6$ и проваливает
positivity. Represented junk quotient и length-two relative block получают
$5/6$, сохраняя одновременно exactness и positivity, но не наследуются.

\begin{theorem}[Endpoint positivity no-go и relative-block fork]
Точный endpoint trace единственно имеет веса $(1,-1)$ и потому не может
быть положительным на полном прямосуммарном curvature space. Положительный
точный маршрут существует только после rank-one проекции на relative
curvature block; такой quotient пока отсутствует в унаследованном calculus.
\end{theorem}

\begin{nogocriterion}[Запрет ручного удаления диагональных блоков]
Нельзя заменить полный curvature trace нормой relative block, просто
отбросив две endpoint-компоненты. Проектор $P_{\rm rel}$ должен следовать
из представленного junk-идеала, степени формы, mapping-cone differential
или эквивалентного унаследованного quotient до вычисления target spectrum.
\end{nogocriterion}

\begin{protocol}\begin{enumerate}
\item endpoint channel rank: \textbf{$2$};
\item unique exact weights: \textbf{$(1,-1)$};
\item positive endpoint trace с exact $Q$: \textbf{невозможен};
\item relative projector rank: \textbf{$1$};
\item projected Gram inertia: \textbf{$(0_-,1_0,1_+)$};
\item inherited relative selector rank: \textbf{$0$};
\item audit rank: \textbf{$6$};
\item прошедших кандидатов: \textbf{$0/12$};
\item лучший positive exact route: \textbf{relative junk quotient, $5/6$};
\item следующий гейт: \textbf{родитель relative-curvature junk quotient}.
\end{enumerate}\end{protocol}

Машинный сертификат сохранён в
\path{s2t_v10_particle_wrinkle_dislocation_callias_equal_charge_twist_m4_h15_r2_superconnection_curvature_trace_selector_candidate_audit_gate_results.json}.